\section{Introduction}

The main contributions of this work are summarized as follows:

Light field imaging captures both spatial and angular information, providing rich geometric cues for depth estimation [1]. Epipolar plane images (EPIs) serve as the primary representation in this domain, relying on the Lambertian reflectance assumption, where pixel intensity remains constant across different viewpoints. This physical assumption enables the formation of clear linear structures in EPIs, allowing conventional algorithms to extract accurate depth maps by calculating line slopes. Consequently, EPI-based architectures have become the standard for light field depth estimation, demonstrating high accuracy in controlled environments with diffuse surfaces.

Real-world environments, however, are predominantly composed of non-Lambertian surfaces, such as specular, translucent, and scattering materials. On these surfaces, the physical assumption of photo-consistency is violated. Specular reflections shift with viewpoints, destroying the linear structures in EPIs, while volumetric scattering eliminates a single well-defined surface depth . Existing EPI-based networks, including EPINet and its directional variants, experience significant performance degradation when applied to such regions. In complex mixed urban scenes or purely non-Lambertian datasets, the mean absolute error (MAE) of conventional EPI models increases considerably. The structural collapse in EPIs leads to slope ambiguity, causing depth estimation errors to exceed acceptable thresholds by a large margin.

The challenge is compounded by a critical data scarcity bottleneck and theoretical limitations in feature extraction. Publicly available non-Lambertian light field datasets contain only a limited number of training scenes (e.g., merely 4 scenes in standard benchmarks), creating a hard barrier for purely data-driven generalization. Previous attempts to mitigate this through attention mechanisms or auxiliary loss functions fail to address the underlying physical decoupling. Under discrete angular resolutions (e.g., $9 \times 9$), traditional frequency-domain analyses like the 2D discrete Fourier transform struggle to extract discriminative material features. This limitation restricts the accurate pixel-level classification of surface reflectance types, preventing existing models from adaptively handling mixed reflectance scenarios.

To address these limitations, we propose a unified dual-mask physical model for non-Lambertian light field depth estimation. Instead of relying on continuous angular frequency analysis, we introduce a three-layer angular signal decomposition physical model. This model decomposes the angular signal $I(u,v)$ into a DC component representing ambient illumination, a linear component $a \cdot u + b \cdot v$ representing disparity and depth, and a residual component $\varepsilon(u,v)$ representing material and reflectance properties. By leveraging geometric angular features, such as angular gradients and correlation coherence, we achieve high-precision material classification under discrete sampling. Based on this physical decomposition, an adaptive geometric dual-mask routing architecture is designed to generate spatial-angular masks, dynamically routing Lambertian and non-Lambertian regions to specialized processing branches.

The main contributions of this work are summarized as follows:
\begin{itemize}
 \item We propose a three-layer angular signal decomposition physical model that decouples illumination, depth, and material information at the geometric level, enabling accurate pixel-level reflectance classification under discrete angular resolutions.
 \item We develop an adaptive geometric dual-mask routing architecture that utilizes the extracted material masks to process Lambertian and non-Lambertian regions through specialized branches, effectively overcoming the failure of single-EPI models on physically violated surfaces.
 \item We integrate a domain-balanced sampling strategy with a weighted random sampler to optimize training stability and cross-domain generalization in complex mixed scenes, validating the framework on multiple benchmarks including HCInew, UrbanLF-Syn, and non-Lambertian datasets.
\end{itemize}
The remainder of this paper is organized as follows. Section II reviews related work on light field depth estimation and non-Lambertian surface modeling. Section III details the proposed three-layer physical decomposition model and the geometric feature extraction process. Section IV presents the adaptive dual-mask routing architecture and the domain-balanced training strategy. Section V provides extensive experimental evaluations and ablation studies. Section VI concludes the paper.

Conventional light field depth estimation methods predominantly rely on epipolar plane image (EPI) (EPI) analysis, operating under the strict Lambertian reflectance assumption. Methodologically, these approaches treat depth estimation as a slope extraction problem in the spatial-angular domain. Technically, they frequently employ spectral analysis techniques, such as 2D discrete Fourier transform (DFT) or structure tensor methods, to identify the dominant orientation of EPI lines . However, this paradigm exhibits two critical limitations. First, spectral methods inherently require high and continuous angular resolution to achieve sufficient frequency domain separation. In practical light field cameras with discrete and sparse angular sampling (e.g., $9 \times 9$ or $5 \times 5$ micro-lens arrays), the frequency resolution degrades significantly, leading to aliasing and inaccurate slope calculation. Second, when applied to non-Lambertian surfaces, the photo-consistency assumption breaks down. Specular highlights shift across viewpoints, creating curved or broken structures in EPIs, while translucent materials cause volumetric scattering that blurs the EPI lines. Consequently, the frequency components of non-Lambertian regions overlap with those of the background or adjacent textures, rendering spectral slope extraction entirely ineffective.

To overcome the handcrafted feature limitations of traditional algorithms, early deep learning methods, typified by EPINet and its directional variants (e.g., EPINet4Dir), formulate depth estimation as a data-driven regression task [2]. While these convolutional neural networks achieve high accuracy on synthetic or purely Lambertian datasets (e.g., HCInew), they suffer from severe performance degradation in real-world mixed scenarios. Methodologically, these models lack explicit physical priors, treating the complex angular variations of non-Lambertian materials as mere noise or relying entirely on the network's implicit capacity to memorize them. Technically, this black-box approach encounters a severe data scarcity bottleneck. Publicly available non-Lambertian light field datasets contain extremely limited training samples (e.g., merely four scenes in standard benchmarks), which prevents the network from learning robust feature representations for specular or scattering regions. A single-branch architecture therefore attempts to process both Lambertian and non-Lambertian pixels using the same set of weights. This induces severe feature conflict during optimization, causing the model to overfit to the majority Lambertian domains while yielding substantial depth errors in minority non-Lambertian and complex urban domains.

Recent advanced approaches attempt to mitigate these issues by incorporating attention mechanisms, multi-scale feature fusion, or defocus-based auxiliary cues . Although these architectural enhancements improve overall metrics on standard benchmarks, they fail to resolve the physical inconsistencies inherent in non-Lambertian reflections. Methodologically, they still process the entire light field through a unified pathway without distinguishing the underlying material properties. Technically, attention modules often mistakenly assign high weights to specular highlights, misinterpreting view-dependent intensity shifts as texture edges or depth discontinuities. Similarly, defocus cues, while complementary to EPI slopes, become unreliable on highly reflective or transparent surfaces where the point spread function deviates from standard geometric optics. Without an explicit material-aware routing mechanism, these networks cannot adaptively direct non-Lambertian regions to specialized processing branches. The reliance on computationally heavy modules also increases inference latency, while the absence of domain-balanced training strategies limits cross-domain generalization when transitioning from controlled laboratory environments to unconstrained mixed scenes.

In summary, existing methods are constrained by their reliance on continuous angular spectral analysis, vulnerability to EPI structural collapse on non-Lambertian surfaces, and the inability to resolve feature conflicts caused by extreme data scarcity in mixed domains. The absence of a physically grounded signal decomposition model and an adaptive material-routing architecture prevents accurate depth recovery in complex real-world environments. Therefore, there is an urgent need for a unified framework that explicitly decouples illumination, depth, and material properties from discrete angular signals, and adaptively routes different reflectance regions to specialized processing branches to achieve robust light field depth estimation.

In this paper, we propose a Unified Dual-Mask Physical Model for non-Lambertian light field depth estimation, which explicitly decouples illumination, depth, and material properties from discrete angular signals. Unlike conventional approaches that rely on continuous spectral analysis or purely data-driven regression, our framework introduces a three-layer angular signal decomposition physical model. This model factorizes the light field angular signal into a direct current (DC) component for ambient illumination, a linear component for disparity, and a residual component for material reflectance. Based on this physical decoupling, we design an adaptive geometric dual-mask routing architecture that directs Lambertian and non-Lambertian regions to specialized processing branches. By leveraging geometric angular features rather than frequency-domain representations, the proposed method effectively handles sparse angular sampling (e.g., $9 \times 9$) and the structural degradation of epipolar plane image (EPI)s (EPIs) caused by specular and translucent surfaces.

\begin{itemize}\item \textbf{Contribution 1: Three-Layer Angular Signal Decomposition Physical Model.} We formulate a physical decomposition framework that factorizes the discrete angular signal into distinct physical components:\end{itemize}
 \begin{equation}
\label{eq:eq1}
I(u,v) = \text{DC} + a \cdot u + b \cdot v + \varepsilon(u,v)
\end{equation}
 where DC represents ambient illumination, the linear term ($a \cdot u + b \cdot v$) encodes disparity, and the residual $\varepsilon(u,v)$ captures material reflectance. By replacing continuous spectral analysis with geometric angular features (e.g., angular gradient and correlation coherence) extracted from the residual component, this model resolves the theoretical infeasibility of material feature extraction under sparse angular sampling. This decoupling provides a robust physical foundation for pixel-level binary classification of Lambertian and non-Lambertian surfaces, achieving a high discriminative power (Cohen's $d = 0.983$).

\begin{itemize}\item \textbf{Contribution 2: Adaptive Geometric Dual-Mask Routing Architecture.} We develop a unified dual-branch network that adaptively routes spatial-angular features based on pixel-level geometric masks. The Lambertian EPI branch utilizes continuous depth priors and EPI line structures, while the non-Lambertian geometric branch exploits angular gradients and X-shape patterns in EPIs to recover depth from structurally degraded regions. This dual-mask mechanism overcomes the failure of single-branch EPI models when physical photo-consistency assumptions are violated by specular or scattering surfaces, ensuring unified and adaptive processing across diverse reflectance types.\end{itemize}

\begin{itemize}\item \textbf{Contribution 3: Domain-Balanced Training Strategy for Mixed Scenes.} We design a domain-balanced sampling strategy integrated with a weighted random sampler to address the severe data scarcity in non-Lambertian scenarios. By maintaining cross-domain exposure balance and employing random tone mapping augmentation, this strategy enhances the training stability and cross-domain generalization of the model in complex, real-world mixed environments, preventing the network from overfitting to the dominant Lambertian domains.\end{itemize}

Detailed experiments on diverse datasets, including HCInew, UrbanLF-Syn, and a data-scarce non-Lambertian dataset, validate the effectiveness of the proposed method. The model achieves an overall Mean Absolute Error (MAE) of less than 0.20, with specific MAEs of $<0.16$ for Lambertian scenes, $<0.22$ for mixed urban scenes, and $<0.25$ for non-Lambertian surfaces, yielding substantial improvements over existing EPI-based baselines and directional EPINet variants.